\section{Busqueda intra-sistema de candidatos a planetas no detectados}

\subsection{Motivacion}

Este modulo traslada la incompletitud topologica global a la escala de sistemas planetarios conocidos. El objetivo no es afirmar detecciones fisicas directas, sino priorizar intervalos orbitales donde podria existir submuestreo intra-sistema.

\subsection{Metodologia}

Cada sistema se ordena por periodo orbital y se analizan gaps adyacentes. Para gaps grandes se generan candidatos sinteticos en escala logaritmica de periodo y semieje mayor, con apoyo adicional de analogos globales y proxies de detectabilidad por transito o velocidad radial.

\subsection{Uso de resultados previos TOI/ATI/sombra}

Cuando existen tablas previas, el modulo incorpora TOI regional, ATI por planeta ancla y sombra observacional por nodo Mapper. Esto convierte la topologia global mas la sombra observacional mas la arquitectura intra-sistema en un ranking de intervalos orbitales prioritarios.

\subsection{Resultados principales}

Los resultados priorizan candidatos topologicos y regiones sugeridas para seguimiento observacional. No se detectan planetas reales ni se afirma un conteo verdadero de planetas ausentes.

\begin{itemize}

\item WASP-132: max score=0.972, n candidatos=6.

\item HD 27894: max score=0.971, n candidatos=4.

\item HD 153557: max score=0.954, n candidatos=4.

\item TOI-4010: max score=0.953, n candidatos=5.

\item KOI-142: max score=0.947, n candidatos=4.

\end{itemize}

\paragraph{Top candidatos}

\begin{itemize}

\item WASP-132: intervalo entre WASP-132 b y WASP-132 d, periodo sugerido=198.097 dias, clase=high.

\item HD 27894: intervalo entre HD 27894 c y HD 27894 d, periodo sugerido=259.749 dias, clase=high.

\item WASP-132: intervalo entre WASP-132 b y WASP-132 d, periodo sugerido=65.416 dias, clase=high.

\item HD 153557: intervalo entre HD 153557 c y HD 153557 d, periodo sugerido=73.943 dias, clase=high.

\item TOI-4010: intervalo entre TOI-4010 d y TOI-4010 e, periodo sugerido=157.110 dias, clase=high.

\end{itemize}

\subsection{Validacion leave-one-out}

La validacion interna retira temporalmente planetas intermedios conocidos y verifica si el modulo vuelve a priorizar candidatos cercanos en periodo orbital.

Se ejecutaron 524 pruebas. El error mediano en log10(P) fue 0.063, con recall dentro de 0.1 dex de 0.544 y dentro de 0.2 dex de 0.761.

\subsection{Limitaciones}

El ranking no modela completitud instrumental ni sustituye analisis de inyeccion-recuperacion. Las estimaciones de masa, radio y detectabilidad son proxies para priorizacion, no inferencias fisicas definitivas.

\subsection{Trabajo futuro}

Los siguientes pasos incluyen incorporar completitud instrumental, seguimiento observacional real y validacion cruzada con catalogos futuros.
